\begin{statement}{E.64}{Proved}{Full hyperelliptic descent and the exact two CM comparison characters}
Let $\Phi_{n,F}$ be the standard braid automorphism associated to
$F=(\lambda,f)\in\widehat{\mathrm{GT}}$.  Write
\[
 \delta_n=\sigma_1\cdots\sigma_{n-1},\qquad z_n=\delta_n^n,
 \qquad
 W_n=\sigma_1\cdots\sigma_{n-2}\sigma_{n-1}^2
                         \sigma_{n-2}\cdots\sigma_1.
 \tag{E.64.1}
\]
For every $g\geq2$, the full profinite hyperelliptic quotient
\[
 \widehat\Upsilon_g=
 \widehat B_{2g+2}/
 \overline{\langle\!\langle
 z_{2g+2},\ W_{2g+2}^2,\ [W_{2g+2},\sigma_i]
 \rangle\!\rangle}
 \tag{E.64.2}
\]
is preserved by $\Phi_{2g+2,F}$.  The induced automorphism fixes its
deck involution and induces the standard action on its full
genus-zero quotient.  In genus two this action has the actual
two-boundary sewing comparison
\[
 \begin{gathered}
 (\widehat B_4\times\widehat{\mathbf Z})/
                   \overline{\langle(z_4,-2)\rangle}
       \ \simeq\ (\widehat\Gamma_2)_{\vec\gamma},
       \qquad t^2=z_4,\\
 (b,d^a)\longmapsto
                (\Phi_{4,F}(b),d^{\lambda a}),
       \qquad t\longmapsto t^\lambda .
 \end{gathered}
 \tag{E.64.3}
\]
Here $\vec\gamma$ is the oriented nonseparating cut curve; both
boundary twists of the cut surface map to $t$.

The two CM subgroup maps are equally explicit.  In
$B_4=\langle a,b,c\rangle$, let
\[
 \begin{gathered}
 q:B_4\longrightarrow B_3=\langle a,b\rangle,
       \qquad q(a)=q(c)=a,\quad q(b)=b,\\
 K=\ker q=\langle\xi_1,\xi_2\rangle,\qquad
 \xi_1=ca^{-1},\quad
 \xi_2=aba c^{-1}b^{-1}a^{-1},\qquad r=ab,
       \quad s=aba .
 \end{gathered}
 \tag{E.64.4}
\]
With $B_3^{\rm mor}=\langle u,v\rangle$, the map
\[
 \theta(u)=ab,\qquad \theta(v)=b^{-1}cba
 \quad\text{identifies}\quad
 \widehat B_3^{\rm mor}\simeq
             \widehat q^{-1}(\overline{\langle r\rangle}).
 \tag{E.64.5}
\]
It maps $(uv)^3$ to $z_4$.  For the quartic subgroup put
\[
 \begin{gathered}
 d=abc,\qquad y=b^{-1}d^2b,\qquad x=(yd)^{-1},\\
 G_4=\langle D,Y\mid Y^2=D^4\rangle,
       \qquad D\mapsto d,\quad Y\mapsto y.
 \end{gathered}
 \tag{E.64.6}
\]
This identifies
$\widehat G_4\simeq\widehat q^{-1}(\overline{\langle s\rangle})$.
The centre maps to $\overline{\langle z_4\rangle}$, and the
central quotient is the specified triangle group
$C_4\amalg C_2$, with generators $(x,y,d)$ of product one.
Both displayed subgroup maps contain exactly the same free
Ferrari kernel $\widehat K$.

Restrict the remaining assertions to $\lambda=1$.  Let $A_F,
\beta_F$ be the unique positive pure frames of E.1 and E.58, lifted
to the free complement $\langle a^2,b^2\rangle$ in $\widehat B_3$.
There is a unique full quartic frame
\[
 H_F=j(\beta_F)k_F,\qquad k_F\in\widehat K,
 \qquad \Phi_{4,F}(d)=H_FdH_F^{-1},
 \qquad q(H_F)=\beta_F,
 \tag{E.64.7}
\]
where $j$ is the first-three-strand section of $q$.  The normalized
actions
\[
 \begin{split}
 \Psi_{3,F}&=\theta^{-1}\operatorname{Inn}(A_F^{-1})
                                    \Phi_{4,F}\theta,\\
 \Psi_{4,F}&=\operatorname{Inn}(H_F^{-1})\Phi_{4,F}
                                  \bigm|_{\widehat G_4}
 \end{split}
 \tag{E.64.8}
\]
are actual group actions.  The first fixes $u$ and $(uv)^3$;
the second fixes $d$ and $z_4$.

In the ordered basis $(\xi_1,\xi_2)$ use
\[
 \begin{gathered}
 U=\begin{pmatrix}1&1\\0&1\end{pmatrix},\quad
 L=\begin{pmatrix}1&0\\-1&1\end{pmatrix},\quad
 R_+=UL=\begin{pmatrix}0&1\\-1&1\end{pmatrix},\quad
 S=ULU=\begin{pmatrix}0&1\\-1&0\end{pmatrix}.
 \end{gathered}
 \tag{E.64.9}
\]
Let $Q$ be this literal action of the first $\widehat B_3$, and
let $M_F$ be the action of $\Phi_{4,F}$ on $\widehat K^{\rm ab}$.
Then
\[
 \begin{gathered}
 M_F=\begin{pmatrix}1&t_F\\0&1\end{pmatrix},\qquad
 M_FL M_F^{-1}=Q(f)^{-1}LQ(f),\qquad t_F\in\widehat{\mathbf Z},\\
 C_3(F)=Q(A_F)^{-1}M_F\in\operatorname{Cent}(R_+),\qquad
 C_2(F)=Q(\beta_F)^{-1}M_F\in\operatorname{Cent}(S).
 \end{gathered}
 \tag{E.64.10}
\]
The first covariance characterizes $M_F$ among matrices commuting
with $U$ and fixing $e_1$, over every commutative coefficient ring.
In particular the signed rigidity shear with that covariance is
literally $M_F$.  Both $C_3$ and $C_2$ are continuous characters
into the corresponding integral norm-one groups:
\[
 \begin{array}{c|c|c}
  &\text{centralizer}&\text{determinant}\\ \hline
 R_+&aI+bR_+&a^2+ab+b^2\\
 S&aI+bS&a^2+b^2.
 \end{array}
 \tag{E.64.11}
\]
On the E.1 CM matrix fibre their reductions at three are exactly
the positive coordinates $T,W$.  The first bit is the sign of
$T$ on $(1,-1)$; the second is zero for $W=\pm I$ and one for
$W=\pm S$.

Finally let $\Omega_{3,F}$ and $\Omega_{4,F}$ be the original
orbifold actions of the same pair on the two displayed triangle
groups, lifted with the central powers specified below, and let
$B_{3,F},B_{4,F}$ be their matrices on the same Ferrari basis.
The comparisons
\[
 \mathcal D_3(F)=C_3(F)B_{3,F}^{-1},\qquad
 \mathcal D_4(F)=C_2(F)B_{4,F}^{-1}
 \tag{E.64.12}
\]
are actual continuous norm-one characters.  On a Galois element
$\sigma$ with $\chi(\sigma)=1$, the opened NTY comparison gives
the stronger full-group formula
\[
 \Psi_{3,\sigma}\Omega_{3,\sigma}^{-1}
       =\operatorname{Inn}(u^{3\rho_K(\sigma)}),\qquad
 \mathcal D_3(\sigma)=(-I)^{\rho_K(\sigma)}.
 \tag{E.64.13}
\]
No value of either comparison character on all of
$\widehat{\mathrm{GT}}$ is assigned in this statement.
\end{statement}
\begin{proof}
We prove the presentation descent and each based subgroup map,
then retain the separate orbifold comparisons.

\paragraph{Full hyperelliptic presentation.}
The exact presentation in
\cite[Theorems 3.2--3.3, pp.6--7]{BCEndoKamada2016}
gives (E.64.2) and its deck involution $W_n$ for $n=2g+2$.
Its positive generators satisfy exactly our Artin relations;
the identity on those generators gives the presentation map from
the source's right-composition convention.  The needed braid
identities are
\[
 \begin{gathered}
 z_{\rm first}=(\sigma_1\cdots\sigma_{n-2})^{n-1},\quad
 z_{\rm last}=(\sigma_2\cdots\sigma_{n-1})^{n-1},\\
 z_{\rm last}=\delta_n z_{\rm first}\delta_n^{-1},\qquad
 W_n=z_n z_{\rm last}^{-1}.
 \end{gathered}
 \tag{E.64.14}
\]
The standard action of
\cite[(1.1)]{BCNakamuraSchneps2000} restricts literally to the
first $n-1$ strands and sends each initial full twist $z_k$ to
$z_k^\lambda$.  Thus, for
$v_n=\Phi_{n,F}(\delta_n)\delta_n^{-1}$,
\[
                 \Phi_{n,F}(W_n)=v_nW_n^\lambda v_n^{-1}.
 \tag{E.64.15}
\]
In the quotient (E.64.2), $W_n$ is a central involution and
$\lambda$ is odd.  Hence every stated relator maps to one.
Applying the same argument to $F^{-1}$ proves equality of the
closed kernels and an actual quotient automorphism.  Killing
$W_n$ is the usual full sphere mapping-class quotient, so its
action is the specified genus-zero action.  The completion
comparison is exact: the finite quotients of a presented discrete
quotient are precisely the finite quotients of its source that
annihilate the displayed relators.

\paragraph{The selected genus-two cut and sewing.}
For the two-boundary torus, the positive three-chain presentation
is
\[
 \Gamma(S_1^2)=B_4\times\langle d\rangle,\qquad
 \delta_1=d,\quad\delta_2=z_4d^{-1}.
 \tag{E.64.16}
\]
This is the boundary version of the Birman--Hilden map;
\cite[Section 6]{BCMargalitWinarski} gives the indicated group,
and the three-chain relation identifies the two boundary factors.
Gluing the boundary components by an annulus sends both twists
to the oriented core twist $t$.  Cutting a homeomorphism that
preserves that orientation shows that the kernel is generated
by $\delta_1\delta_2^{-1}$: its ambiguity is exactly twisting
on the gluing annulus.  This gives the discrete version of
(E.64.3).

For the intrinsic completed stabilizer, use the actual stabilizer
exact sequence in \cite[Theorem 2.3]{BCBoggi2026v6} and its
genus-at-most-two congruence comparison in Proposition 2.6(i).
The completed group on the left of (E.64.3), after killing $t$,
is $\widehat B_4/\langle z_4\rangle=
\widehat{P\Gamma}_{1,2}$.  Its central $t$ is an injected
$\widehat{\mathbf Z}$: the map
$(b,a)\mapsto\operatorname{exp}(b)+6a$ annihilates $(z_4,-2)$
and sends $t$ to six.  The natural map to the actual stabilizer
is the identity on this inertia and on the indicated cut-surface
quotient.  The two exact sequences prove the claimed isomorphism.

In the genus-two chain $T_1,\ldots,T_5$, take $t=T_5$ and
$B_4=\langle T_1,T_2,T_3\rangle$.  The three-chain relation is
$(T_1T_2T_3)^4=T_5^2$.  Also
$z_5=(T_1T_2T_3T_4)^5$ is the deck involution: on the sphere
it is the twist around five of the six marked points, hence is
trivial, and on integral homology it is $-I$.  Therefore
$f(z_5,T_5^2)=1$, since its two arguments commute and $f$ is a
commutator word.  The standard formula gives
$\Phi_{6,F}(T_5)=T_5^\lambda$, and its restriction to the first
three generators is $\Phi_{4,F}$.  This proves the actual
sewing square, including the two central boundary images.

\paragraph{The Mordell subgroup, before completion.}
The Ferrari free kernel and its Nielsen action are the literal
maps of \cite[(2.9), Lemma 4.6]{DE_NakamuraTsunogaiYasuda2010}:
\[
 \begin{array}{c|cc}
  &\xi_1&\xi_2\\ \hline
 a&\xi_1&\xi_2\xi_1\\
 b&\xi_1\xi_2^{-1}&\xi_2.
 \end{array}
 \tag{E.64.17}
\]
The following full braid words retain the central signs in the
Mordell triangle:
\[
 \begin{gathered}
 \xi_+=ab,\qquad
 \xi_{+2}=b^{-1}a^{-1}b^{-1}a^{-1}b^{-1}c^{-1},\qquad
 \xi_{+3}=cbab,\\
 \theta(uv)=\xi_{+3},\quad
 \theta(uvu)=\xi_{+2}^{-1},\quad
 \theta((uv)^3)=z_4,\\
 \xi_+\xi_{+2}\xi_{+3}=1,\qquad
 \xi_{+2}^2=z_4^{-1},\quad\xi_{+3}^3=z_4.
 \end{gathered}
 \tag{E.64.18}
\]
They prove the Artin relation for $\theta$.  Both folded
generators equal $r$.  In the source put
\[
 k_1=(uv)u^{-2},\qquad
 k_2=u^4(uv)(uv)^{-3}.
 \tag{E.64.19}
\]
Their images are exactly $\xi_1,\xi_2$.  The derived subgroup
of $B_3$ is free of rank two: its image in $C_2*C_3$ is the
index-six free subgroup, and the full centre has nonzero
exponent sum.  The map from that free group onto $K$, also
free of rank two, is an isomorphism by Hopficity.  Thus the two
displayed words form a basis, and $\theta$ is an isomorphism on
the free kernels and infinite cyclic quotients.  This proves
the discrete version of (E.64.5).  In particular the centre
map is $z_3^{\rm mor}\mapsto z_4$, not a discarded central power.

\paragraph{The quartic subgroup, before completion.}
Direct braid multiplication gives
\[
 \begin{gathered}
 y^2=d^4,\qquad q(d)=s,\quad q(y)=s^2,\quad q(x)=s^{-3},\\
 k_0=yd^{-2}=\xi_2^{-1},\qquad
 k_1=dyd^{-3}=\xi_1^{-1},\\
 d\xi_1d^{-1}=\xi_2^{-1},\qquad
 d\xi_2d^{-1}=\xi_1.
 \end{gathered}
 \tag{E.64.20}
\]
There is a complete presentation verification: define
$\alpha(k_0)=k_1$, $\alpha(k_1)=k_0^{-1}$ on the free group.
The mutually inverse maps
\[
 \begin{gathered}
 G_4\longrightarrow\langle k_0,k_1\rangle\rtimes_\alpha\mathbf Z,
       \qquad D\mapsto(1,1),\quad Y\mapsto(k_0,2),\\
 k_0\mapsto YD^{-2},\qquad k_1\mapsto DYD^{-3},\qquad
                         1\in\mathbf Z\mapsto D
 \end{gathered}
 \tag{E.64.21}
\]
are homomorphisms because $\alpha^2(k_0)=k_0^{-1}$ and
$\alpha^4=1$.  Equation (E.64.20) identifies both free kernels,
and $D\mapsto s$ identifies their cyclic quotients.  This proves
the discrete quartic assertion.  Its centre is
$\langle D^4\rangle$, and killing it gives $C_4*C_2$.
The kernel of this free product onto $C_4$, with $D\mapsto1$
and $Y\mapsto2$, is exactly $K$.  Its cusp and centre relation is
\[
 \begin{split}
 x&=\xi_1d^{-3},\\
 x^4z_4^3&=[\xi_1,\xi_2^{-1}]
          =\xi_2^{-1}[\xi_1,\xi_2]^{-1}\xi_2.
 \end{split}
 \tag{E.64.22}
\]
The algebraic cover $Y^2=X^3-X$, $t=1/X^2$, has precisely
these ramification indices $(4,2,4)$, with the origin removed
over $t=0$ and deck generator $(X,Y)\mapsto(-X,iY)$.
The displayed triangle generators have deck values $(1,2,1)$.
The exact braid basis, rather than an unspecified path to that
algebraic cover, is used in all matrices below.

\paragraph{Completion of the two actual inclusions.}
The split extension $B_4=K\rtimes B_3$ induces the full
profinite topology on $K$.  Indeed, every characteristic finite
quotient of the finitely generated free group $K$ has a finite
automorphism group, so the action factors through a finite
quotient of $B_3$ and supplies the corresponding finite
semidirect quotient of $B_4$.  Such characteristic subgroups
are cofinal among the finite-index subgroups of $K$.
The section splits the completed quotient, whence
$\widehat B_4=\widehat K\rtimes\widehat B_3$.
The closures of $\langle r\rangle$ and $\langle s\rangle$ are
each $\widehat{\mathbf Z}$, since the exponent-sum maps take
them to two and three respectively.  The same construction
completes $B_3^{\rm mor}=K\rtimes\mathbf Z$ and (E.64.21).
It proves both stated completed subgroup isomorphisms.

\paragraph{Full accessory frames and their cocycles.}
The E.1 pure frames give the equalities modulo the first
three-strand centre.  Their exponent sums agree upstairs, and
that centre has exponent sum six.  Hence the exact lifted
equalities are
\[
 \Phi_{3,F}(r)=A_FrA_F^{-1},\qquad
 \Phi_{3,F}(s)=\beta_Fs\beta_F^{-1}.
 \tag{E.64.23}
\]
The Ferrari fold and section satisfy
$q\Phi_{4,F}=\Phi_{3,F}q$ and
$\Phi_{4,F}j=j\Phi_{3,F}$ on their literal generators.
Thus normalization by $A_F$ gives the first map in (E.64.8).

For the quartic map, E.39 supplies a full conjugator $H$ for
$d$.  The centralizer of $s$ in $\widehat B_3$ is
$\overline{\langle s\rangle}$: after quotienting by
$z_3=s^2$, the factor-centralizer theorem
\cite[Theorem B', Section 4, p.159]{BCHerfortRibes1985}
gives $C_{C_2\amalg C_3}(\bar s)=\langle\bar s\rangle$.
Consequently $q(H)=\beta_Fs^a$ for a unique
$a\in\widehat{\mathbf Z}$.  Replace $H$ by $Hd^{-a}$.
This preserves its conjugacy equation and gives (E.64.7).

This frame is unique.  The same theorem gives
$C_{C_4\amalg C_2}(\bar d)=\langle\bar d\rangle$.
A member of $\widehat K$ centralizing $d$ therefore has
trivial image in the central quotient, since
$\widehat K\cap\langle\bar d\rangle=1$ there.  It is then
in $\overline{\langle d^4\rangle}\cap\widehat K=1$, as
$q(d^4)=s^4$ has infinite profinite order.  Thus
$C_{\widehat K}(d)=1$.

The frames are continuous by the compact conjugation maps and
their uniqueness.  That uniqueness gives the exact identities
\[
 \begin{gathered}
 A_{FG}=\Phi_{3,F}(A_G)A_F,\qquad
 \beta_{FG}=\Phi_{3,F}(\beta_G)\beta_F,\\
 H_{FG}=\Phi_{4,F}(H_G)H_F.
 \end{gathered}
 \tag{E.64.24}
\]
For the last one, the right side has precisely the required
projection $\beta_{FG}$.  Equations (E.64.24) prove that both
normalized maps in (E.64.8) are group actions.  They fix the
stated generators and centres literally.

\paragraph{Integral homology and the rigidity identification.}
Equation (E.64.17) gives $Q(a)=U$, $Q(b)=L$.  In the standard
formula, $e=f(z_3,c^2)$ lies in $\widehat K$, because
$q(e)=f(z_3,a^2)=1$.  Both its arguments commute with $a$,
so $e$ does also.  Therefore
$\Phi_{4,F}(\xi_1)=e^{-1}\xi_1e$, and the first column of
$M_F$ is $e_1$.  Since $\Phi_{4,F}(a)=a$, covariance gives
$M_FU=UM_F$.  This proves the asserted unipotent shear.
Covariance with
$\Phi_{3,F}(b)=f(a^2,b^2)^{-1}bf(a^2,b^2)$ gives the
first identity in (E.64.10).

Here is the promised comparison over every coefficient ring.
For $A=\left(\begin{smallmatrix}a&b\\c&d\end{smallmatrix}\right)$,
\[
 AU-UA=\begin{pmatrix}-c&a-d\\0&c\end{pmatrix},\qquad
 AL-LA=\begin{pmatrix}-b&0\\a-d&b\end{pmatrix}.
 \tag{E.64.25}
\]
Thus the joint centralizer consists of scalars, over any
commutative ring.  The quotient of two matrices satisfying
the same covariance commutes with $U,L$; fixing $e_1$ makes
that scalar one.  This proves the exact equality with the
signed rigidity shear, including over $\mathbf Z_\ell$ and
$\widehat{\mathbf Z}$.

For the finite E.61 convention put $t=8\rho$ and
\[
 F_\rho=\begin{pmatrix}1&-t\\-t&1+t^2\end{pmatrix},\qquad
 V=U^t,\qquad R_0=-R_+.
 \tag{E.64.26}
\]
Direct integer multiplication gives
\[
 \begin{gathered}
 VLV^{-1}=F_\rho^{-1}LF_\rho,\qquad
 VSV^{-1}=SF_\rho,\qquad
 VR_0V^{-1}=F_\rho^{-1}R_0.
 \end{gathered}
 \tag{E.64.27}
\]
Consequently $t_F=8\rho_F\pmod {12}$, with precisely this
orientation.  On the E.1 matrix fibre $F_\rho=I$, so
$M_F=I\pmod3$.

The first normalized action fixes $u$, whose matrix on the
free kernel is $R_+$, and the second fixes $d$, whose matrix
is $S$.  Their homology matrices are exactly $C_3,C_2$:
the factor $k_F$ of (E.64.7) is an inner automorphism of
$\widehat K$ and hence has trivial abelianization action.
The frame cocycles and covariance prove multiplicativity.
Solving the four entry equations gives (E.64.11) over every
commutative ring; the determinants are one.  At three the
norm-one groups have the six and four elements used in E.1.
Their matrices are the positive $T,W$ by (E.64.23) and
$M_F=I\pmod3$, with the exact two bit maps stated above.

\paragraph{The original orbifold maps, with their central lifts.}
In the Mordell source put
$x=u$, $y=(uvu)^{-1}$, $z=uv$.  Then $xyz=1$ and
$y^2=z_3^{-1}$, $z^3=z_3$.  In the quartic source use the
three elements $(x,y,z)=(x,y,d)$ of (E.64.6); here
$y^2=z^4=z_4$.  In either source the original puncture-loop
action has the same actual formula
\[
 \begin{split}
 \Omega_F(x)&=x^\lambda,\\
 \Omega_F(y)&=f(y,x)y^\lambda f(x,y),\\
 \Omega_F(z)&=x^m f(z,x)z^\lambda f(x,z)x^{-m},
                         \qquad m=(\lambda-1)/2.
 \end{split}
 \tag{E.64.28}
\]
Modulo the stated centre, its product relation is the gentle
equation, and the finite-order generators retain their orders.
The product discrepancy upstairs is central.  In the Mordell
case the exponent sum sends that centre to six; in the
quartic case the map $(z,y,x)\mapsto(1,2,-3)$ sends it to four.
The discrepancy has value zero, hence is one.  The central
power relations map to their $\lambda$ powers by conjugation,
so (E.64.28) gives actual lifts.  The inverse pair and the
same central detection prove invertibility and the composition
law.  This proves the orbifold maps as separate homomorphisms;
it makes no unstated identification with the standard
three-strand braid action.

At multiplier one, the Mordell map fixes $x=u$ and the
quartic image of $d$ is conjugate to $d$ by $f(d,x)\in K$.
Their matrices therefore commute respectively with $R_+$
and $S$.  Their determinants are one.  For example the exact
cusp relations are
\[
 [k_1,k_2]u^6=z_3^{\rm mor},\qquad
 x^4z_4^3=[\xi_1,\xi_2^{-1}],
 \tag{E.64.29}
\]
so both maps fix a primitive commutator in the second
lower-central quotient.  Their homology representations
$B_{3,F},B_{4,F}$ have the same bases as $C_3,C_2$.
All four matrices lie in the respective commutative
centralizers.  Their quotients are therefore the actual
characters (E.64.12).  For $B_{3,F}$, this is exactly the
triangle Tate action of
\cite[Proposition 4.4]{DE_NakamuraTsunogaiYasuda2010}.

\paragraph{The evaluated Galois comparison and its scope.}
The specified path in
\cite[Proposition 4.11, Section 4.12, Proposition 4.13]
 {DE_NakamuraTsunogaiYasuda2010}, at multiplier one, has
\[
 u_\sigma=h_\sigma(a^2,b^2)^{-1}(ab)^{3\rho_K(\sigma)}
           =A_\sigma z_3^{\rho_K(\sigma)}.
 \tag{E.64.30}
\]
Here $h_\sigma=A_\sigma^{-1}$ is the equianharmonic frame;
the harmonic frame denoted $g_\sigma$ in that source is
$\beta_\sigma^{-1}$.  The source's full braid-generator
equations give
$\Phi_{4,\sigma}\theta=
\operatorname{Inn}(u_\sigma)\theta\Omega_{3,\sigma}$.
Since $z_3=(ab)^3=\theta(u^3)$, normalizing by $A_\sigma$
gives (E.64.13) at the full group level.  On homology the
first-three-strand centre acts by $R_+^3=-I$.
The total centre instead acts trivially, as is also recorded
by the exact boundary word
\[
                   [\xi_1,\xi_2]=z_4z_3^{-2}.
 \tag{E.64.31}
\]
These explicit centre maps account for the sign in the
Galois comparison.  Neither (E.64.15) nor the selected
sewing square replaces the unevaluated full-GT values of
(E.64.12) by (E.64.30).

\paragraph{Retained source audit and exact certificates.}
The source audit preserves the actual capping map from
$B_{2g+1}/\langle z_{2g+1}^2\rangle$ as a branch-point
stabilizer of index $2g+2$.  It also preserves the identity
and $\operatorname{Inn}(\tau_{\gamma_i})$ as two different
based automorphisms agreeing on the oriented curve
stabilizer in the setting of Boggi v6 Lemma 6.16.  The
present proof uses the full presentation and the explicit
sewing maps above; it does not require either printed
identification.  The exact audit and the tested rejected
double-centre alternative remain in the research files.

There are two finite replay programs, sharing prefix
\texttt{bc\_}: \texttt{hyperelliptic\_sewing.py} and
\texttt{cm\_norm\_characters.py}.  The former verifies 45
literal braid identities with the faithful Artin action
and the separate integral genus-two matrix comparisons.
The latter verifies 14 quartic braid identities, the
independent free-by-cyclic presentation, the signed
covariances, and both positive CM tables.  Its checks of
centralizers over $\mathbf Z/m$, $2\leq m\leq12$, supplement
the coefficient-ring proof (E.64.25).  Each program retains
its source hashes and complete \texttt{.json} and
\texttt{.out} records.  They assign no full-GT value to the
comparison characters beyond the maps proved here.
\end{proof}
